\chapter{Hereditary Hemochromatosis Genetic Test}

\section*{What Is This Test?}
This blood test checks for common inherited changes in the HFE gene, especially C282Y and H63D. Some of these changes can make the body absorb too much iron over many years. Extra iron may damage the liver, heart, pancreas, or joints.

\section*{Alternative Names}
\begin{itemize}
  \item HFE Gene Test
  \item Hemochromatosis Mutation Test
  \item C282Y and H63D Test
  \item Hereditary Iron Overload Genetic Test
\end{itemize}
\section*{How It Is Performed}
A lab extracts DNA from a blood sample and checks for the common HFE changes. The result is read together with iron tests, especially transferrin saturation and ferritin. The gene test alone cannot show how much iron is stored in the body.

\section*{Why It Is Ordered}
\begin{itemize}
  \item To evaluate persistently high transferrin saturation or ferritin when iron overload is suspected
  \item To investigate unexplained liver disease, diabetes, cardiomyopathy, arthropathy, or bronze skin pigmentation when the iron studies look compatible
  \item To test adult first-degree relatives of a person with confirmed HFE-related hemochromatosis
  \item To clarify a suspected inherited cause of iron overload
\end{itemize}

\section*{Interpretation and Limitations}
A person with two C282Y changes has an increased chance of HFE-related iron overload. However, not everyone with this result develops illness. One C282Y and one H63D change usually carries a lower risk. A negative result does not rule out other inherited forms of iron overload or other causes of high ferritin. The clinician also reviews iron tests, liver health, symptoms, and family history.
\section*{Specimen and Preparation}
Fasting is not needed for the gene test. Other iron tests may have different instructions. Tell the clinician about alcohol use, inflammation, liver disease, blood transfusions, iron supplements, and relatives with iron overload. These can change the iron results.

\section*{Risks and Safety}
The blood draw may cause brief pain or bruising. The gene result does not measure current iron levels or show whether an organ is damaged. Genetic counseling can help when the result may affect relatives.

\section*{Understanding Results and Follow-up}
A positive result is usually followed by iron tests and, when needed, checks of the liver or other organs. A negative common-variant test does not rule out every cause of iron overload. Adult first-degree relatives may be offered testing after counseling. Treatment is based on iron levels and health findings, not the gene result alone.

\section*{Regulatory Status and Authorization Date}
This chapter describes a test category, not a specific commercial product. FDA, CDSCO, EU IVDR, and MHRA approval or registration dates therefore cannot be assigned without the assay manufacturer, product name, instrument, and intended use. For laboratory-developed tests, an individual agency approval date may not exist. See the Preface for the interpretation of regulatory status.

\clearpage
